\title{Controlling Text-to-Image Diffusion by Orthogonal Finetuning}

\begin{abstract}
State-of-the-art text-to-image diffusion models have demonstrated remarkable proficiency in producing lifelike images based on textual descriptions. A fundamental open question is how to direct or steer these models effectively for a variety of downstream applications. We propose a systematic finetuning strategy, Orthogonal Finetuning (OFT), designed to adapt text-to-image diffusion models for new tasks. Distinct from prior approaches, OFT provides theoretical guarantees for preserving hyperspherical energy—a measure reflecting the pairwise relationships of neurons confined to the unit hypersphere. We identify the preservation of this property as critical for maintaining the semantic fidelity of generated images in text-to-image models. To further enhance the robustness of finetuning, we introduce Constrained Orthogonal Finetuning (COFT), which adds an explicit constraint on the hypersphere's radius. Our investigation focuses on two primary finetuning tasks for text-to-image models: subject-driven generation, aimed at creating individualized subject images from limited examples paired with text prompts, and controllable generation, which empowers models to incorporate auxiliary control signals. Extensive empirical studies reveal that the OFT framework achieves superior performance in terms of both image generation quality and the speed of convergence, surpassing existing techniques.
\end{abstract}

\section{Introduction}

The latest generation of text-to-image diffusion models~\cite{saharia2022photorealistic,ramesh2022hierarchical,rombach2022high} excel in producing high-quality, photorealistic images guided by textual instructions. However, these models sometimes struggle with ambiguity or lack the precision required for granular, accurate manipulation of the image content based on text alone. In response to this limitation, this paper concentrates on two key directions within text-to-image generation:

\begin{itemize}[leftmargin=*,nosep]

    \item \textbf{Subject-driven generation}~\cite{ruiz2023dreambooth}: When only a small set of images depicting a particular subject is available, this task focuses on generating novel images of that subject in varied scenarios, steered by accompanying text prompts.
    \item \textbf{Controllable generation}~\cite{zhang2023adding,mou2023t2i}: Here, the model is tasked with image generation conditioned on both textual prompts and supplementary control inputs (such as canny edge maps or segmentation masks), adhering to the specified control signals.
\end{itemize}

At their core, both tasks involve optimizing the finetuning of text-to-image diffusion models in such a way that the quality of the original pretrained model’s generative abilities is maintained. The principal criteria for a robust finetuning method can be summarized as follows: (1) \emph{training efficiency}, characterized by a reduction in trainable parameters and required training epochs, and (2) \emph{generalizability preservation}, which entails retaining the capacity for high-fidelity, diverse image synthesis. Conventional finetuning approaches typically proceed either by making incremental updates to neuron weights using a small learning rate (\emph{e.g.}, \cite{ruiz2023dreambooth}), or by introducing a modest re-parameterized module with new neuron weights (\emph{e.g.}, \cite{hulora2022,zhang2023adding}). However, despite their straightforward nature, these strategies do not offer formal guarantees for retaining the generative competence acquired during pretraining. Furthermore, there is insufficient theoretical guidance for crafting optimal finetuning procedures or selecting appropriate hyperparameters like training duration. This challenge largely stems from the absence of reliable metrics for assessing how well the pretrained model’s generative skills are preserved after finetuning. Current techniques often implicitly posit that a smaller Euclidean distance between the finetuned and original models suggests better retention of pretrained capabilities, leading to the prevalent use of extremely low learning rates. Although this relationship may sometimes hold, we contend that Euclidean distance alone fails to adequately reflect the semantic continuity between models. Consequently, leveraging more structurally-informed metrics to evaluate the disparity between the finetuned and pretrained models could substantially enhance both the fidelity of pretraining retention and the stability of the finetuning process.

Drawing upon prior empirical findings that hyperspherical similarity effectively captures semantic relationships~\cite{liu2017deep,liu2018decoupled,chen2020angular}, we leverage the concept of hyperspherical energy~\cite{liu2018learning} to represent the inter-neuron relational patterns within a layer. Hyperspherical energy aggregates the hyperspherical similarity (\emph{e.g.}, cosine similarity) across all neuron pairs in a layer, thereby quantifying the extent to which neuron representations are evenly distributed on the unit hypersphere~\cite{liu2021learning}. We conjecture that, for optimal finetuning, the hyperspherical energy should exhibit minimal deviation from that of the original pretrained network. A straightforward strategy would involve introducing a regularization term to ensure that hyperspherical energy is conserved during finetuning; however, this approach does not inherently guarantee close alignment between the energies before and after finetuning. To address this limitation, we utilize a key invariance of hyperspherical energy: applying an identical orthogonal transformation to every neuron preserves their pairwise hyperspherical similarities. Building on this insight, we introduce Orthogonal Finetuning~(OFT), an approach designed to adapt large text-to-image diffusion models to specific tasks by learning a shared orthogonal transformation per layer, ensuring that hyperspherical energy remains unchanged. The main objective is to maintain pairwise angles among neurons by optimizing a layer-wise orthogonal mapping, which can also be seen as reorienting the neurons within a new canonical coordinate framework. Further considering that reduced Euclidean distance between the finetuned and pretrained models correlates with better retention of the pretrained features, we extend our method with a variant termed Constrained Orthogonal Finetuning~(COFT). COFT imposes an additional constraint, ensuring that the parameters of the finetuned model reside on a hypersphere of fixed radius centered at the pretrained neurons.

The underlying rationale for employing orthogonal transformations during neuron finetuning draws partial inspiration from the 2D Fourier transform, which decomposes an image into its magnitude and phase spectra. Notably, the phase spectrum—encoding the angular relationships between the input and basis elements—retains the essential semantic information. For instance, one can reconstruct the main content of an image using its phase spectrum combined with a randomized magnitude spectrum, without significant loss of meaning. This observation implies that adjusting neuron orientations is fundamental for inducing semantic changes in generated images—a central aim in both subject-driven and controllable image generation. Nevertheless, allowing neuron directions to change with excessive flexibility tends to compromise the original generative capacity developed during pretraining. To restrict such flexibility, we advocate maintaining the angles between every pair of neurons, resting on the premise that these angular relationships encode important representational knowledge within neural networks. Guided by this reasoning, we propose applying an orthogonal transformation shared across each layer, thereby ensuring that the hyperspherical energy remains invariant.

Additionally, our work is motivated by orthogonal over-parameterized training~\cite{liu2021orthogonal}, where classification neural networks are trained from a random initialization by means of layer-wise orthogonal transformations. This approach leverages the fact that randomly initialized networks generally exhibit low expected hyperspherical energy; the goal in \cite{liu2021orthogonal} is to keep hyperspherical energy low throughout training, as reduced energy correlates with better generalization in classification tasks~\cite{liu2018learning,lin2020regularizing}. Their results demonstrate that orthogonal transformations can provide enough expressiveness to facilitate the training of generalizable classifiers. In contrast, our focus is on finetuning text-to-image diffusion models to enhance both controllability and generative capability for downstream tasks. Two key differences distinguish OFT from the method in \cite{liu2021orthogonal}. Firstly, whereas \cite{liu2021orthogonal} is oriented toward minimizing hyperspherical energy, our objective with OFT is to preserve the hyperspherical energy of the pretrained model, safeguarding the structural integrity established during pretraining—a critical concern, as reducing hyperspherical energy in diffusion model finetuning could impair essential semantic representations. Secondly, OFT is explicitly formulated to minimize divergence from the pretrained model, yielding a constrained version of the transformation. In contrast, the method in \cite{liu2021orthogonal} does not introduce such constraints. Finetuning fundamentally involves balancing adaptability against preservation, and we posit that our OFT paradigm provides an effective solution to this challenge. Below, we summarize our contributions:

\begin{itemize}[leftmargin=*,nosep]

    \item We introduce Orthogonal Finetuning, a new approach for finetuning text-to-image diffusion models with enhanced controllability. To bolster the robustness of the finetuning process, we further present a constrained variant that bounds the angular shift from the original pretrained model.
    \item In contrast to conventional finetuning techniques, OFT uniquely guarantees preservation of hyperspherical energy during model adaptation—an aspect we empirically identify as crucial for maintaining the generative semantic integrity inherent to the pretrained architecture.
    \item We utilize OFT for two distinct settings: subject-driven generation and controllable generation. Through an extensive suite of experiments, we establish that our method achieves substantial improvements over previous strategies with respect to generated image fidelity, speed of convergence, and reliability during finetuning. In addition, OFT is markedly more sample efficient, showing effective convergence using far fewer training samples and training epochs.
    \item For the controllable generation context, we propose a novel metric—control consistency—to rigorously assess how well controllability is preserved. This metric operates by extracting the control cue from each synthesized output and quantifying its correspondence to the original control instruction. Our findings affirm that OFT delivers superior controllability.
\end{itemize}

\section{Related Work}

\textbf{Text-to-image diffusion models}. The landscape of text-to-image synthesis has experienced remarkable advancements~\cite{saharia2022photorealistic,ramesh2022hierarchical,rombach2022high,gu2022vector,nichol2022glide}, propelled chiefly by the evolution of diffusion-based generative paradigms~\cite{ho2020denoising,song2021score,song2021denoising,dhariwal2021diffusion} alongside advances in vision-language representation frameworks~\cite{su2020vlbert,lu2019vilbert,radford2021learning,wang2021simvlm,li2022blip,li2023blip,alayrac2022flamingo,schuhmann2022laion}. In terms of architecture, GLIDE~\cite{nichol2022glide} and Imagen~\cite{saharia2022photorealistic} operate directly within pixel space: GLIDE achieves joint training of a text encoder and diffusion prior with paired text-image inputs, whereas Imagen employs a fixed, pretrained text encoder. Conversely, Stable Diffusion~\cite{rombach2022high} and DALL-E2~\cite{ramesh2022hierarchical} develop generative models over latent spaces; Stable Diffusion leverages VQ-GAN~\cite{esser2021taming} for a visual codebook underpinning its latent representations, while DALL-E2 utilizes CLIP~\cite{radford2021learning} to map both images and captions into a shared embedding space. Beyond diffusion-based systems, the domain also investigates generative adversarial networks~\cite{reed2016generative,zhang2017stackgan,xu2018attngan,li2019controllable} and autoregressive architectures~\cite{ramesh2021zero,ding2021cogview,wu2022nuwa,yuscaling} for similar synthesis objectives. OFT represents a flexible, model-agnostic finetuning mechanism that is compatible with any text-to-image diffusion framework.

\textbf{Subject-driven generation}. Existing solutions for preserving subject characteristics in generation tasks include user-supplied masks as supplementary conditioning signals~\cite{avrahami2022blended,nichol2022glide}. Alternative strategies, such as inversion-based approaches~\cite{choi2021ilvr,dhariwal2021diffusion,ramesh2022hierarchical,gal2022image}, adjust contextual content while keeping the depicted subject unchanged, and \cite{hertz2022prompt} enable both localized and global editing absent of explicit masks. Nevertheless, these methods are often inadequate for accurately retaining intricate identity features of the target subject. Pivotal Tuning~\cite{roich2022pivotal} addresses this by adaptively refining the generator in the vicinity of an inverted latent representation with identity-preserving regularization. Likewise, \cite{nitzan2022mystyle} constructs a tailored face prior using a set of personal facial images. \cite{casanova2021instance} facilitates instance-level variation synthesis, though instance-specific detail may be compromised. DreamBooth~\cite{ruiz2023dreambooth}, by incorporating customized tokens and a handful of reference images, finetunes the diffusion model using both reconstruction and class-prior preservation losses. Our proposed OFT approach builds on the DreamBooth framework, replacing conventional finetuning at low learning rates with updates constrained to orthogonal

\textbf{Controllable generation}. Image-to-image translation is fundamentally a controllable generation problem. Historically, many approaches have employed conditional generative adversarial networks for this task~\cite{isola2017image,zhu2017toward,wang2018high,park2019semantic,choi2018stargan}. More recently, diffusion models have also been explored for image-to-image translation~\cite{saharia2022palette,voynov2022sketch,wang2022pretraining}. Notably, ControlNet~\cite{zhang2023adding} enhances a pretrained diffusion model by finetuning it with supplemental control signals, leading to highly effective controllable generation. In parallel, T2I-Adapter~\cite{mou2023t2i} achieves increased controllability in generation by similarly finetuning diffusion models. Building upon these approaches~\cite{zhang2023adding,mou2023t2i}, we introduce OFT for finetuning pretrained diffusion models, achieving superior controllability with reduced requirements for training data and finetuning parameters. Importantly, OFT maintains the same computational efficiency at inference, incurring no additional costs during test-time.

\textbf{Model finetuning}. The practice of adapting large-scale pretrained models to downstream tasks has gained considerable traction in recent years~\cite{devlin2018bert,brown2020language,he2022masked}. Within this context, adaptation techniques (\emph{e.g.}, \cite{hulora2022,houlsby2019parameter,pfeiffer2020adapterfusion}) have been widely examined, particularly in natural language processing. OFT is most closely related to LoRA~\cite{hulora2022}, which applies low-rank decomposition to model weight adjustments during finetuning. Distinctly, OFT leverages layer-wise shared orthogonal transformations to multiplicatively update neuron weights, ensuring the preservation of pairwise neuron angles within a layer and thereby improving stability.

\section{Orthogonal Finetuning}

\subsection{Why Does Orthogonal Transformation Make Sense?}

We begin by motivating the use of orthogonal transformations for finetuning text-to-image diffusion models. This question can be divided into two parts: (1) understanding the importance of adjusting the orientation (angle) of neurons, and (2) justifying the choice of orthogonal transformations as the mechanism for modifying neuron directions during finetuning.

For the first question, our approach is informed by the empirical findings of \cite{liu2018decoupled,chen2020angular}, which note that the angular discrepancy between features is an effective indicator of the semantic gap. SphereNet~\cite{liu2017deep} demonstrates that constraining all neurons to lie on a unit hypersphere during training enables a neural network to retain a comparable modeling capacity and, in many cases, improves generalization. This observation points to the neuron directions as the primary carrier of significant data characteristics. To underscore the significance of neuron angles, we devise a simple experiment, illustrated in Figure~\ref{fig:toy_ae}: we train a conventional convolutional autoencoder with a dataset of flower images. During training, the standard inner product forms the feature map ($z$ refers to an individual output of the convolution filter $\bm{w}$ applied to input $\bm{x}$ within a local region). In evaluation, three alternative methods are applied for feature map computation: (a) replicating the training’s inner product, (b) isolating magnitude information, and (c) leveraging only angular information. Results, depicted in Figure~\ref{fig:toy_ae}, reveal that recovering images from angular data alone is nearly perfect, while magnitude alone yields no discernible information. Notably, the cosine-based (c) activation is not present during training, which relies solely on inner products. These results reinforce that neuron direction predominantly encodes the semantic properties of the inputs. Consequently, adjusting neuron directions appears more promising for modulating semantic aspects of images.

Addressing the second question, a straightforward method to update neuron directions is to apply a collective rotation or reflection across all neurons within a layer, a strategy that invokes orthogonal transformation. While more granular angular transformations—allowing each neuron a unique rotation—could offer greater flexibility, we have observed that orthogonal transformations offer an optimal balance between adaptability and regularity. Furthermore, \cite{liu2021orthogonal} provides evidence for the expressive power of orthogonal transformations in training neural networks. To further substantiate this perspective, we present an experiment evaluating the regularization effect imparted by the orthogonality condition. Utilizing the ControlNet~\cite{zhang2023adding} framework, we conduct a controlled generation experiment; the findings, shown in Figure~\ref{fig:ortho}, indicate that our canonical OFT maintains stability and achieves high-precision control by training epoch 20. By contrast, when the orthogonality constraint is removed, OFT fails to produce plausible images or exert any meaningful control. This experimental outcome highlights the critical role of orthogonality in ensuring the efficacy of OFT.

\subsection{General Framework}

Traditional finetuning approaches often employ gradient descent with a low learning rate to adjust all or select layers of a model. By using a small learning rate, these methods implicitly restrict updates, ensuring that the finetuned weights do not stray far from those of the pretrained model. In effect, conventional finetuning amounts to optimizing the model parameters while indirectly minimizing the distance $\thickmuskip=2mu \medmuskip=2mu \| \bm{M} - \bm{M}^0\|$, where $\bm{M}$ and $\bm{M}^0$ represent the finetuned and pretrained weight matrices, respectively. This soft regularization is intuitively appealing, but may still permit too much freedom when adapting large-scale models. To mitigate this, LoRA introduces a stricter low-rank constraint, enforcing $\thickmuskip=2mu \medmuskip=2mu \text{rank}(\bm{M} - \bm{M}^0)=r'$, where $r'$ is a small predefined integer. In contrast, OFT instead constrains the similarity between pairs of neurons by ensuring $\thickmuskip=2mu \medmuskip=2mu \| \text{HE}(\bm{M})-\text{HE}(\bm{M}^0) \|=0$, with $\text{HE}(\cdot)$ denoting the hyperspherical energy associated with a weight matrix. 

For illustration, consider a fully connected layer parameterized by $\thickmuskip=2mu \medmuskip=2mu \bm{W}=\{\bm{w}_1,\ldots,\bm{w}_n\}\in\mathbb{R}^{d\times n}$, where each $\bm{w}_i\in\mathbb{R}^{d}$ refers to the $i$-th neuron's weight vector (with pretrained weights $\bm{W}^0$). Given an input $\thickmuskip=2mu \medmuskip=2mu \bm{x}\in\mathbb{R}^d$, the output $\thickmuskip=2mu \medmuskip=2mu \bm{z}\in\mathbb{R}^n$ is determined by $\thickmuskip=2mu \medmuskip=2mu \bm{z}=\bm{W}^\top\bm{x}$. In the OFT paradigm, the learning objective becomes the minimization of the difference in hyperspherical energy between the updated and original models:

\begin{equation}\label{eq:oft_principle}
\footnotesize
\min_{\bm{W}} \left\| \text{HE}(\bm{W})-\text{HE}(\bm{W}^0)\right\|~~\Leftrightarrow~~\min_{\bm{W}} \bigg{\|} \sum_{i\neq j} \|\hat{\bm{w}}_i-\hat{\bm{w}}_j\|^{-1}-\sum_{i\neq j} \|\hat{\bm{w}}^0_i-\hat{\bm{w}}^0_j\|^{-1}\bigg{\|}
\end{equation}

Here, $\thickmuskip=2mu \medmuskip=2mu \hat{\bm{w}}_i:=\bm{w}_i/\|\bm{w}_i\|$ indicates the normalized version of the $i$-th neuron's weights, and hyperspherical energy for a layer is defined by $\thickmuskip=2mu \medmuskip=2mu \text{HE}(\bm{W}):= \sum_{i\neq j} \|\hat{\bm{w}}_i-\hat{\bm{w}}_j\|^{-1}$. The optimal value of Eq.~\eqref{eq:oft_principle} is clearly zero, a minimum attained whenever $\bm{W}$ and $\bm{W}^0$ differ only by a rotation or reflection; that is, when $\thickmuskip=2mu \medmuskip=2mu \bm{W}=\bm{R}\bm{W}^0$ for some orthogonal matrix $\thickmuskip=2mu \medmuskip=2mu \bm{R}\in\mathbb{R}^{d\times d}$ (the determinant indicating rotation for $1$ or reflection for $-1$). This forms the fundamental principle behind OFT, where the neural network is finetuned by training orthogonal, layer-wise shared matrices

\begin{equation}\label{eq:oft_general}
    \footnotesize
    \bm{z}=\bm{W}^\top\bm{x}=(\bm{R}\cdot\bm{W}^0)^\top\bm{x},~~~\text{s.t.}~\bm{R}^\top\bm{R}=\bm{R}\bm{R}^\top=\bm{I}
\end{equation}

In this formulation, $\bm{W}$ corresponds to the weight matrix fine-tuned using OFT, while $\bm{I}$ represents the identity matrix. The OFT process is depicted in Figure~\ref{fig:block}. Drawing inspiration from the zero initialization employed in LoRA, it is essential for OFT to begin fine-tuning from the pretrained model’s parameters $\bm{W}^0$. To ensure this, the orthogonal matrix $\bm{R}$ is initially set as the identity matrix, which causes the starting point for fine-tuning to coincide with the pretrained weights. Maintaining the orthogonality of $\bm{R}$ can be accomplished with differentiable orthogonalization approaches, as outlined by \cite{liu2021orthogonal,lezcano2019cheap}. We will elaborate on strategies for enforcing orthogonality that are also computationally scalable.

\subsection{Efficient Orthogonal Parameterization}\label{sect:eff_ortho}

Although the Gram-Schmidt procedure can provide differentiable orthogonalization, it typically incurs high computational cost in practice~\cite{liu2021orthogonal}. To enhance computational efficiency, we employ the Cayley transform to parameterize the orthogonal matrix. More precisely, the orthogonal matrix is formed as $\thickmuskip=2mu \medmuskip=2mu \bm{R}=(\bm{I}+\bm{Q})(I-\bm{Q})^{-1}$, where $\bm{Q}$ denotes a skew-symmetric matrix such that $\thickmuskip=2mu \medmuskip=2mu \bm{Q}=-\bm{Q}^\top$. This approach, while efficient, introduces a limitation—Cayley parameterization only yields orthogonal matrices with determinant $1$, restricting $\bm{R}$ to the special orthogonal group. However, our empirical results indicate that this restriction does not compromise practical effectiveness. Despite the Cayley transform, for high-dimensional $\bm{R}$ the parameter count may still be prohibitive. To mitigate this issue, we design $\bm{R}$ as a block-diagonal matrix with $r$ constituent blocks, resulting in the structure:

\begin{equation}
    \footnotesize
    \bm{R}=\text{diag}(\bm{R}_1,\bm{R}_2,\cdots,\bm{R}_r)=\begin{bmatrix}
    \bm{R}_1\in \text{O}(\frac{d}{r}) &  &\\
    &\ddots & \\
    & & \bm{R}_r\in \text{O}(\frac{d}{r})
    \end{bmatrix}\in \text{O}(d)
\end{equation}

where $\text{O}(d)$ refers to the orthogonal group of degree $d$, with $\thickmuskip=2mu \medmuskip=2mu \bm{R}\in\mathbb{R}^{d\times d}$ and each $\thickmuskip=2mu \medmuskip=2mu \bm{R}_i\in\mathbb{R}^{d/r\times d/r},\forall i$ representing orthogonal matrices. When $\thickmuskip=2mu \medmuskip=2mu r=1$, the block-diagonal structure of the orthogonal matrix simplifies to the conventional unstructured orthogonal matrix. For a standard $d\times d$ orthogonal matrix, the number of free parameters is given by $\thickmuskip=2mu \medmuskip=2mu d(d-1)/2$, which corresponds to a parameter complexity of $\mathcal{O}(d^2)$. In contrast, for an $r$-block diagonal orthogonal matrix, the parameter count is reduced to $\thickmuskip=2mu \medmuskip=2mu d(d/r-1)/2$, yielding a complexity of $\thickmuskip=2mu \medmuskip=2mu \mathcal{O}(d^2/r)$. Further parameter savings can be achieved by reusing the same block matrix across all blocks, i.e., $\thickmuskip=2mu \medmuskip=2mu \bm{R}_i = \bm{R}_j, \forall i\neq j$, lowering the complexity to $\mathcal{O}(d^2/r^2)$. Importantly, these parameter-efficient schemes maintain the orthogonality of $\bm{R}$, ensuring that hyperspherical energy is fully retained.

Let us analyze the relative parameter efficiency of OFT in comparison with LoRA. For LoRA using a rank $r'$ decomposition, its trainable parameter count is given by $\thickmuskip=2mu \medmuskip=2mu r'(d+n)$. Assuming $r$ and $r'$ scale with the model dimension $d$ (for example, setting $\thickmuskip=2mu \medmuskip=2mu r = r' = \alpha d$, with $\thickmuskip=2mu \medmuskip=2mu 0<\alpha \leq 1$ as a constant), LoRA's parameter complexity scales as $\mathcal{O}(d^2+dn)$, whereas OFT can reach $\mathcal{O}(d)$ complexity. For concrete comparison, if we use a weight matrix sized $\thickmuskip=2mu \medmuskip=2mu 128\times 128$, LoRA requires $2,048$ trainable parameters for $\thickmuskip=2mu \medmuskip=2mu r'=8$, but OFT only needs $960$ trainable parameters when $\thickmuskip=2mu \medmuskip=2mu r=8$, without applying block sharing.

\subsection{Constrained Orthogonal Finetuning}

We can further restrict the expressiveness of the standard OFT by enforcing that the finetuned model stays in close proximity to the original pretrained model. Concretely, COFT applies the following forward computation:

\begin{equation}\label{eq:coft_general}
    \footnotesize
    \bm{z}=\bm{W}^\top\bm{x}=(\bm{R}\cdot\bm{W}^0)^\top\bm{x},~~~\text{s.t.}~\bm{R}^\top\bm{R}=\bm{R}\bm{R}^\top=\bm{I},~\norm{\bm{R}-\bm{I}}\leq \epsilon
\end{equation}

This enforces both the orthogonality of $\bm{R}$ and bounds its deviation from the identity matrix by $\epsilon$. While orthogonality can be readily imposed through the Cayley parameterization discussed in Section~\ref{sect:eff_ortho}, introducing the $\epsilon$-deviation constraint into a Cayley-parameterized matrix is less straightforward. To gain a better understanding of the Cayley transform in this context, we approximate $\thickmuskip=2mu \medmuskip=2mu \bm{R}=(\bm{I}+\bm{Q})(I-\bm{Q})^{-1}$ using the Neumann series, obtaining $\thickmuskip=2mu \medmuskip=2mu

Given that the series converges under the operator norm, the constraint $\thickmuskip=2mu \medmuskip=2mu\|\bm{R}-\bm{I}\|\leq\epsilon$ can equivalently be imposed within the Cayley transform framework, translating to $\thickmuskip=2mu \medmuskip=2mu \|\bm{Q}-\bm{0}\|\leq\epsilon'$ for some $\epsilon'$ (distinct from $\epsilon$), where $\bm{0}$ represents a zero matrix. This alternative constraint on $\bm{Q}$ is straightforward to implement via projected gradient descent. For initialization that recovers the identity in $\bm{R}$, we simply initialize $\bm{Q}$ as the zero matrix. Conceptually, COFT introduces two explicit restrictions: a minimal difference in hyperspherical energy and a controlled departure from the original pretrained model. While the latter is often an implicit assumption in many finetuning strategies, COFT formalizes it as an explicit condition. In practice, even though OFT achieves strong results, the explicit constraint on the deviation enhances finetuning stability further in COFT, as evidenced in Figure~\ref{fig:eps_coft}. This figure illustrates how the value of $\epsilon$ influences COFT performance: larger $\epsilon$ values allow the finetuned model to mimic the OFT-trained model closely, whereas smaller $\epsilon$ values cause the COFT-adapted model to behave more like the original pretrained text-to-image diffusion model.

\subsection{Re-scaled Orthogonal Finetuning}

We introduce a straightforward modification to standard OFT by jointly optimizing a scaling coefficient for each neuron, inspired by the fact that the normalization inherent in hyperspherical energy makes it insensitive to such scaling. The updated forward computation is: $\thickmuskip=2mu \medmuskip=2mu\bm{z}=(\bm{R}\bm{W}^0\bm{D})^\top\bm{x}$\footnote{\textbf{Errata}: In the NeurIPS camera ready version, the re-scaled OFT forward computation was incorrectly given as $\thickmuskip=2mu \medmuskip=2mu\bm{z}=(\bm{D}\bm{R}\bm{W}^0)^\top\bm{x}$, although the implementation was correct and results in Appendix~\ref{app:rescaled} are unaffected.} where $\thickmuskip=2mu \medmuskip=2mu\bm{D}=\text{diag}(s_1,\cdots,s_n)\in\mathbb{R}^{n\times n}$ is a learnable diagonal matrix whose positive entries $s_1,\ldots,s_n$ adjust the scale of individual neurons. Unlike the original OFT in Eq.~\eqref{eq:oft_general}, which only tunes $\bm{R}$, this approach learns both the orthogonal transformation $\bm{R}$ and the scaling factors in $\bm{D}$. This adjustment marginally increases parameter count while granting OFT greater expressive capacity. For experimental consistency and to highlight the isolated impact of orthogonal transformations, we retain the original OFT setup in main results, but additional findings in Appendix~\ref{app:rescaled} indicate that the re-scaled variant offers further improvements.

\section{Intriguing Insights and Discussions}

\textbf{OFT is broadly applicable across neural network architectures}. In principle, OFT can be integrated into any neural network. For example, while LoRA is typically deployed on attention weights for Transformers~\cite{hulora2022}, our experiments focus OFT on attention weights to ensure comparability with LoRA. Furthermore, OFT is naturally compatible with convolutional layers, especially due to the structural properties of the block-diagonal $\bm{R}$, a feature not shared by LoRA. Specifically, by setting the number of blocks equal to the number of input channels, each block operates on a separate neuron channel, analogous to employing depth-wise convolution kernels~\cite{chollet2017xception}. Alternatively, sharing the block matrices across all channels implements an orthogonal transformation identically across neuron channels. Additional preliminary experiments on applying OFT to convolutional layer finetuning are reported in Appendix~\ref{app:convolution}.

\textbf{Connection to LoRA}. LoRA achieves preservation of the information contained within the pretrained weight matrix by introducing an additional low-rank matrix, thereby mitigating drastic shifts. In contrast, OFT regulates the transformation applied to the pretrained weight matrix so that it remains orthogonal (and thus full-rank), ensuring the integrity of the pretrained knowledge is not compromised. The forward pass of OFT can be reformulated as $\thickmuskip=2mu \medmuskip=2mu \bm{z}=(\bm{R}\bm{W}^0)^\top\bm{x}=(\bm{W}^0+(\bm{R}-\bm{I})\bm{W}^0)^\top\bm{x}$, where $\thickmuskip=2mu \medmuskip=2mu (\bm{R}-\bm{I})\bm{W}^0$ plays a similar role to the low-rank adaptation introduced in LoRA. Since $\bm{W}^0$ is normally full-rank, OFT achieves a low-rank update in cases where $\thickmuskip=2mu \medmuskip=2mu \bm{R}-\bm{I}$ is itself low-rank. Analogous to LoRA’s use of a rank parameter $r'$, OFT utilizes a diagonal block parameter $r$ to limit the trainable parameters. Fundamentally, LoRA and OFT embody distinct strategies for parameter efficiency: LoRA leverages a low-rank architecture for reducing learnable parameters, whereas OFT achieves efficiency by exploiting structured sparsity (specifically, block-diagonal orthogonality).

\textbf{Why OFT converges faster}. Figure~\ref{fig:toy_ae} demonstrates that the optimal method for altering semantic content involves changing neuron orientations—a mechanism that OFT is explicitly designed to perform. Additionally, OFT may be interpreted as refining neurons within a smooth hyperspherical manifold, which provides a more favorable optimization landscape. This claim is further supported by empirical findings in \cite{liu2021orthogonal}.

\textbf{Why not minimize hyperspherical energy}. It is important to distinguish our approach from that of \cite{liu2021orthogonal}, as we do not target minimization of hyperspherical energy. For classification tasks, the goal is to maintain distinct, non-redundant neurons. Achieving minimum hyperspherical energy would distribute neurons evenly across the hypersphere, but this criterion does not serve the finetuning process, since it could overwrite or eliminate important information from pretraining.

\textbf{Balancing flexibility and regularity in finetuning}. Our findings reveal an intrinsic balance between flexibility and regularity during finetuning. Conventional finetuning approaches, while providing maximum flexibility, often result in poor model robustness and are prone to catastrophic collapse. OFT, although straightforward in design, strikes a suitable compromise by retaining the angles between neuron pairs, achieving stability alongside adaptability. In addition, employing a block-diagonal structure can be interpreted as imposing stricter regularization constraints on the orthogonal transformation matrix.

\textbf{Inference efficiency remains unchanged}. Distinct from ControlNet, our OFT methodology incurs zero additional computational cost at inference time for the adapted model. During inference, the transformation can be executed by directly left-multiplying the learned orthogonal matrix $\bm{R}$ with the original weight matrix $\bm{W}^0$, leading to an equivalent parameter matrix $\thickmuskip=2mu \medmuskip=2mu \bm{W}=\bm{R}\bm{W}^0$. Consequently, the inference time matches that of the base pretrained model.

\section{Experiments and Results}

\textbf{Experimental protocols}. We conduct our empirical evaluation using Stable Diffusion v1.5~\cite{rombach2022high} as the foundational text-to-image generator. To guarantee objective comparisons, samples are drawn randomly from each experimental setting. For tasks centered on subject-driven generation, our experimental procedure is based largely on DreamBooth~\cite{ruiz2023dreambooth}. For controllable generation, we adhere to methodologies established in ControlNet~\cite{zhang2023adding} and T2I-Adapter~\cite{mou2023t2i}. To ensure consistency with LoRA, both OFT and COFT modifications are confined to layers utilized by LoRA. Comprehensive details and further results are available in Appendix~\ref{app:settings}.

\subsection{Subject-driven Generation}

\textbf{Experimental setup}. DreamBooth~\cite{ruiz2023dreambooth} and LoRA~\cite{hulora2022} serve as our primary baseline models. All evaluated methods employ the loss function defined by DreamBooth. For both DreamBooth and LoRA, we align with the hyperparameter recommendations and implementation details of their original publications. Supplemental results can be found in Appendix~\ref{app:settings},~\ref{app:coft_oft},~\ref{app:more_qualitative}, and~\ref{app:failure}.

\textbf{Finetuning stability and convergence}. We begin by analyzing the stability and convergence rate during finetuning for DreamBooth, LoRA, OFT, and COFT, as depicted in Figure~\ref{fig:teaser} and Figure~\ref{fig:db_convergence}. It is apparent that both OFT and COFT are capable of achieving stable finetuning when applied to the diffusion model. After completing 400 iterations, both DreamBooth and OFT variants manage effective control of the output, whereas LoRA does not maintain the subject's identity. Progressing to 2000 iterations, DreamBooth begins to exhibit mode collapse with degraded outputs, and LoRA fails to generate a yellow shirt, instead producing yellow fur. Conversely, OFT and COFT continue to demonstrate reliable and steady control over generated images. These observations support the assertion that the OFT framework enables both rapid and robust convergence for subject-driven tasks. It is also noteworthy that the low sensitivity of OFT and COFT to the specific number of finetuning iterations mitigates the need for iteration-dependent parameter tuning across different subjects. Thus, it is feasible to preset a generous iteration count for both OFT and COFT without fine-grained adjustment. For COFT, an appropriate choice of $\epsilon$ further simplifies the selection of both learning rate and iteration count, reducing the complexity of hyperparameter tuning.

\textbf{Quantitative comparison}. Adopting the evaluation protocol from \cite{ruiz2023dreambooth}, we present a quantitative analysis addressing subject fidelity (using DINO~\cite{caron2021emerging}, CLIP-I~\cite{radford2021learning}), text prompt fidelity (CLIP-T~\cite{radford2021learning}), and diversity of samples (LPIPS~\cite{zhang2018unreasonable}). For subject fidelity, CLIP-I measures the mean cosine similarity of CLIP embeddings between the generated outputs and reference images, while DINO uses the same metric but replaces CLIP with ViT S/16 DINO embeddings. CLIP-T quantifies the alignment between text prompts and generated images through the average cosine similarity of CLIP embeddings. Sample diversity is computed as the mean LPIPS cosine similarity across generated images corresponding to the same subject and prompt. Results in Table~\ref{table:dreambooth_final} indicate that COFT and OFT surpass DreamBooth and LoRA significantly on both DINO and CLIP-I metrics, while maintaining equal or slightly improved performance on prompt alignment and diversity metrics. For all experiments, we perform finetuning of an identical pretrained model across 30 distinct random seeds to account for variance. The empirical findings confirm that the OFT framework offers enhanced convergence stability and more reliably superior overall performance.

\textbf{Qualitative comparison}. To provide clearer insights into the strengths of OFT, we present a selection of randomly sampled instances of subject-driven generation in Figure~\ref{fig:dreambooth_final}. In order to ensure a fair evaluation, all samples are produced from the identical finetuned model across different methods, without tailoring any text prompt for a specific output. Each method utilizes the checkpoint that yields the highest validation performance according to the CLIP metric. Examining the illustrations in Figure~\ref{fig:dreambooth_final}, it is evident that both OFT and COFT achieve robust semantic retention of the subject, whereas LoRA frequently struggles with subject identity retention (\emph{e.g.}, in the bowl case, LoRA fails to maintain subject consistency). Simultaneously, OFT and COFT demonstrate superior ability to follow textual instructions, in contrast to DreamBooth, which, while capable of preserving subject identity, often falls short in conforming to the given prompts (\emph{e.g.}, see the first row under the bowl example). These observations collectively suggest that the OFT framework excels in both controllability and preservation of the subject. Additionally, OFT exhibits a remarkable stability across various iteration counts; its effectiveness does not deteriorate with more iterations, a property not shared by either DreamBooth or LoRA. Further qualitative illustrations are available in Appendix~\ref{app:more_qualitative}. We also present results from a human evaluation in Appendix~\ref{app:human}, which further supports the advantages of OFT.

\subsection{Controllable Generation}

\textbf{Settings}. For our experiments, we benchmark against ControlNet~\cite{zhang2023adding}, T2I-Adapter~\cite{mou2023t2i}, and LoRA~\cite{hulora2022}. The evaluation focuses on three complex controllable generation scenarios described in the main paper: transforming Canny edge maps to images~(C2I) using the COCO dataset~\cite{lin2014microsoft}, converting segmentation maps to images~(S2I) with the ADE20K dataset~\cite{zhou2017scene}, and mapping landmarks to faces~(L2F) on the CelebA-HQ dataset~\cite{karras2018progressive,xia2021tedigan}. All approaches finetune Stable Diffusion~(SD) v1.5 on each respective dataset over 20 training epochs. Additional results are provided in Appendix~\ref{app:more_qualitative},~\ref{app:more_control_task},~\ref{app:failure}.

\textbf{Convergence}. We assess the rate of convergence for ControlNet, T2I-Adapter, LoRA, and COFT within the context of the L2F task, employing both numerical and visual analyses. For our quantitative metric, the mean $\ell_2$ distance is calculated between the predicted face landmarks and their corresponding control landmarks. Figure~\ref{fig:face_convergence} presents the error of face landmarks as a function of training epochs for each finetuned model. Notably, COFT and OFT exhibit much quicker convergence relative to their counterparts. For instance, LoRA requires 20 epochs to attain the level of accuracy that OFT achieves by epoch 8. It is important to highlight that OFT and COFT maintain a parameter count on par with LoRA—substantially fewer than that of ControlNet—while attaining faster convergence compared to established approaches. Further evidence of OFT's rapid convergence is visible in Figure~\ref{fig:teaser}, where the rightmost example demonstrates OFT's high data efficiency: with merely 5\% of the ADE20K dataset, OFT converges effectively, unlike ControlNet and LoRA. For qualitative comparisons, attention is given to OFT, COFT, and ControlNet, since ControlNet’s landmark error is closest to that of our method. As illustrated in Figure~\ref{fig:face_convergence_qual}, OFT and COFT display consistent and steady convergence, with generated face poses progressively matching the control landmarks. By contrast, ControlNet exhibits a delayed, abrupt enhancement in controllability post-epoch 8, mirroring the sudden convergence described in \cite{zhang2023adding}. The reliable and gradual convergence of OFT significantly eases its practical application, offering smoother and more foreseeable training behavior, thereby simplifying the search for optimal OFT hyperparameters.

\textbf{Quantitative comparison}. To assess controllable generation, we propose a control consistency metric. This involves extracting the control signal from the generated output and measuring its correspondence with the initial input control. Specifically, for the C2I task, the evaluation uses IoU and F1 score, whereas for the S2I task, we employ mean IoU along with both mean and overall accuracy. The L2F task is assessed by calculating the average $\ell_2$ distance between predicted and actual control landmarks. Appendix~\ref{app:settings} provides further specifics regarding these consistency measures. Each method under comparison is evaluated using its optimal hyperparameter configuration. The results in Table~\ref{table:control_gen} indicate that both OFT and COFT deliver superior and more reliable control than competing approaches. Notably, adapter-based models such as T2I-Adapter and ControlNet display slower convergence and inferior final metrics. LoRA outperforms ControlNet on S2I, yet underperforms on C2I and L2F; in addition, we note that LoRA’s performance cap remains relatively modest despite extensive hyperparameter tuning. By contrast, our empirical observations reveal that OFT’s performance continues to improve as we increase the number of trainable parameters, indicating that its potential is not fully reached. Lastly, our systematic evaluation framework for controllable generation marks a key innovation of this work, as it enables precise measurement of finetuned models’ control capacities, subject to the underlying accuracy of the utilized segmentation/detection models.

\textbf{Qualitative comparison}. In addition to quantitative results, we conduct a qualitative analysis of OFT and COFT in relation to leading contemporary techniques such as ControlNet, T2I-Adapter, and LoRA. As shown by the randomly sampled outputs in Figure~\ref{fig:control_final}, OFT and COFT generate visually convincing, high-quality images while ensuring precise control over the outputs. For the S2I task, LoRA is unable to produce images consistent with the provided segmentation maps, whereas ControlNet, OFT, and COFT are all capable of adhering to the specified structure. Notably, OFT and COFT generate more realistic and detail-rich images with superior geometric coherence compared to ControlNet, despite utilizing significantly fewer parameters. In the C2I task, OFT and COFT can synthesize semantically coherent images using only coarse Canny edges, whereas T2I-Adapter and LoRA exhibit inferior performance. Regarding the L2F task, our method demonstrates highly accurate pose control even for faces with complex orientations. Across these three control scenarios, both OFT and COFT consistently outperform previous state-of-the-art approaches, highlighting the effectiveness of the OFT framework for controlled image synthesis. For broader comparison, we include additional qualitative samples for each control task in Appendix~\ref{app:control_exp}, and further showcase the versatility of OFT on additional tasks (such as dense pose-to-human body, sketch-to-image, and depth-to-image) in Appendix~\ref{app:more_control_task}.

\section{Concluding Remarks and Open Problems}

This work is inspired by the insight that angular relationships among neural activations play a pivotal role in encoding visual semantics. We introduce a straightforward yet powerful fine-tuning approach—orthogonal finetuning—aimed at enhancing control in text-to-image diffusion models. Our method focuses on two key areas: subject-driven generation and controllable image synthesis. When compared to prior techniques, OFT achieves improved controllability and greater stability during fine-tuning, all while requiring fewer tunable parameters. Significantly, OFT maintains computational efficiency by avoiding extra inference costs, making it suitable for deployment.

OFT gives rise to several compelling avenues for further exploration. First, ensuring orthogonality relies on Cayley parametrization, which necessitates computing a matrix inverse—a process that can constrain the scalability of OFT. Although employing a block diagonal parametrization provides a partial remedy, efficiently and differentiably inverting matrices remains a nontrivial problem. Second, OFT possesses promising capabilities for compositionality: orthogonal matrices generated from multiple OFT finetuning procedures can be multiplied, and the result is still orthogonal. Investigating whether such a set of orthogonal matrices reliably preserves information from all downstream tasks is a promising research direction. Lastly, the parameter efficiency of OFT hinges on the block diagonal structure, a constraint that may introduce extra bias and restrict expressiveness. Identifying strategies to further enhance parameter efficiency without sacrificing flexibility or introducing bias stands as an important unresolved challenge.

\end{document}